% ------------------------------------------------------------------------------
% Chapter 2
% ------------------------------------------------------------------------------
\chapter{Test Plan}
\label{cha:test_plan}

\section{Features Under Test}

The features to be tested are divided into tree groups.
The idea is to test both normal execution and edge cases for each group.

\begin{itemize}
    \item \textbf{[C] Control dependent (Consumer/Producer)}
    Reset behavior, Input ready/valid handshake/stall, Output ready/valid handshake/stall.
    \item \textbf{[D] Data dependent (Input/output)}
    edge cases (primes/coprimes, multiples, min value, max value), early exits (zero operands, equal operands), and correct output.
    \item \textbf{[S] State machine}
    Transitions between \texttt{IDLE}, \texttt{RUN}, and \texttt{DONE}, along with their cycle timings.
\end{itemize}

\section{Tool for Testing}

To test these features we will use the following tools:
\begin{itemize}
    \item \textbf{[F] FORMAL}: here we will verify the logic from inside the dut.
    The analysis will touch the expected control output and the FSM states, hold, and transitions.
    \item \textbf{[U] UVM}: here we will verify the logic as a black box through stimulus.
    Using a golden model we will verify the correctness of the data.
    \item \textbf{[H] HARDWARE}: here we will verify the hardware timing and correct timings.
    Assertions based approach is used in a formal model.
\end{itemize}

\section{Traceability Matrix}

Here we introduce all the requirements extracted from the specification.
It is divided in the previously presented features and sub-features accompanied by the tool used to verify it.
\begin{table}[H]
    \centering
    \begin{tabularx}{\textwidth}{|c|c|l|X|c|}
        \toprule
        \textbf{Req ID} & \textbf{Feature}  & \textbf{Sub-Feature}  & \textbf{Description}                                                                      & \textbf{Tool} \\ \midrule
        REQ\_000        & C                 & Reset                 & Correct init and graceful recover from reset mid-transaction                              & F/U           \\ \midrule
        REQ\_001        & C                 & Input accept          & Producer \texttt{in\_valid} is high and DUT \texttt{in\_ready} is high accepts input      & F/U           \\ \midrule
        REQ\_002        & C                 & Input stall           & DUT \texttt{in\_ready} is low and producer \texttt{in\_valid} is high                     & H/F           \\ \midrule
        REQ\_003        & C                 & Output accept         & Consumer \texttt{out\_ready} is high and DUT \texttt{out\_valid} is high accepts output   & H/F           \\ \midrule
        REQ\_004        & C                 & Output stall          & Consumer \texttt{out\_ready} is low and DUT \texttt{out\_valid} is high                   & H/F           \\ \midrule
        REQ\_005        & C                 & Output stability      & \texttt{gcd\_out} is stable while \texttt{out\_valid} is high                             & H/F           \\ \bottomrule \toprule
        REQ\_006        & D                 & All operands          & Correct \texttt{gcd\_out} for general operands                                            & U             \\ \midrule
        REQ\_007        & D                 & Edge cases            & Correct \texttt{gcd\_out} for primes/coprimes, multiples, and min/max values              & U             \\ \midrule
        REQ\_008        & D                 & Zero operands         & Correct \texttt{gcd\_out} for $a=0$ and/or $b=0$, early exit                              & U             \\ \midrule
        REQ\_009        & D                 & Equal operands        & Correct \texttt{gcd\_out} for $a=b$, early exit                                           & U             \\ \bottomrule \toprule
        REQ\_010        & S                 & States                & FSM state is either \texttt{IDLE}, \texttt{RUN}, or \texttt{DONE}                         & F             \\ \midrule
        REQ\_011        & S                 & State hold            & FSM holds its state strictly for the necessary duration                                   & F             \\ \midrule
        REQ\_012        & S                 & Transitions           & FSM transitions following the specification                                               & F             \\ \midrule
        REQ\_013        & S                 & Assigned signal       & FSM determines output signals (\texttt{in\_valid}, \texttt{out\_ready})                   & F             \\ \bottomrule
    \end{tabularx}
\end{table}

This second table is meant to associate each tool and relative test to the requirement it tries to verify.
A description explains what the test is actually looking for.
\begin{table}[H]
    \centering
    \begin{tabularx}{\textwidth}{|c|p{3.9cm}|X|p{1.7cm}|}
        \toprule
        \textbf{Tool}   & \textbf{Test name}                                                                & \textbf{Description}                                                                                                      & \textbf{REQ ID}                           \\ \midrule
        F               & \texttt{inputs}                                                                   & Assume a compliant producer that holds \texttt{in\_valid} and data stable during stalls                                   & REQ\_002                                  \\ \midrule
        F               & \texttt{\_reset}                                                                  & Prove DUT initializes to \texttt{IDLE} and clears all registers on reset                                                  & REQ\_000                                  \\ \midrule
        F               & \texttt{\_state}                                                                  & Prove FSM strictly operates within valid, mutually exclusive states                                                       & REQ\_010                                  \\ \midrule
        F               & \texttt{idle\_idle}, \texttt{run\_run}, \texttt{done\_done}                       & Prove FSM correctly holds current state when transition conditions are not met, including output stalls                   & REQ\_004, REQ\_011                        \\ \midrule
        F               & \texttt{idle\_done}, \texttt{idle\_run}, \texttt{run\_done}, \texttt{done\_idle}  & Prove FSM accurately transitions between states based on handshakes and data path evaluations                             & REQ\_001, REQ\_003, REQ\_012              \\ \midrule
        F               & \texttt{\_in\_ready}, \texttt{\_out\_valid}                                       & Prove control output signals are strictly tied to their respective FSM states                                             & REQ\_013                                  \\ \midrule
        F               & \texttt{outputs}                                                                  & Prove \texttt{gcd\_out} stability while \texttt{out\_valid} is asserted                                                   & REQ\_005                                  \\ \bottomrule \toprule
        H               & \texttt{chk\_out\_valid\_no\_drop}                                                & Check that the DUT does not drop \texttt{out\_valid} while waiting for the consumer \texttt{out\_ready} to assert         & REQ\_004                                  \\ \midrule
        H               & \texttt{chk\_out\_data\_stable}                                                   & Check that the DUT holds \texttt{gcd\_out} stable while \texttt{out\_valid} is asserted                                   & REQ\_005                                  \\ \midrule
        H               & \texttt{chk\_reset}                                                               & Check that in case of reset the DUT clear all of its outputs                                                              & REQ\_000                                  \\ \midrule
        H               & \texttt{chk\_early\_exit}                                                         & Check that the DUT skips to the \texttt{DONE} state in case of early exit                                                 & REQ\_008, REQ\_009                        \\ \bottomrule \toprule
        U               & \texttt{scb.write\_rst}                                                           & Verify that the golden model resets its internal state mid-transaction                                                    & REQ\_000                                  \\ \midrule
        U               & \texttt{scb.write\_in}, \texttt{scb.write\_out}                                   & Verify that the golden model output is equal to the DUT \texttt{gcd\_out}                                                 & REQ\_006, REQ\_007, REQ\_008, REQ\_009    \\ \bottomrule
    \end{tabularx}
\end{table}
        

\section{Coverage}

The target coverage is just a measure of what we would like to tackle in this test plan for each area of the dut.
In practice the objectives are a complete (100\%) coverage of the functional description and possibly a high (\textless 90\%) code coverage.
This excludes toggle as it is dependant of the width of the data.

\section{Stimulus}

The stimulus is divided into multiple tests that tackle different scenarios.
\begin{itemize}
    \item \textbf{gcd\_bringup\_test}: Test the basic functionality of the verification system as a whole
    \item \textbf{gcd\_det\_test}: Test the edge cases and tries to obtain the maximum coverage with the minimal number of tests
    \item \textbf{gcd\_rnd\_test}: Test the DUT with random instructions to find edge cases that have not been tested otherwise
\end{itemize}
